\subsection{阿廷定理}

定理 2（阿廷）.——设 N 是一个域，\(\Gamma\) 是 N 的一个自同构群，K 是 \(\Gamma\) 的不变量域。设 V 是 N 的一个向量子 K-空间，它在 K 上有限维。则 V 到 N 中的每个 K-线性映射 u 都是 \(\Gamma\) 的元素在 V 上的限制的、系数在 N 中的线性组合。

设 \(u\) 是 \(\mathbf{V}\) 到 \(\mathbf{N}\) 中的一个 K-线性映射，设 \(\mathrm{V}_{(\mathrm{N})} = \mathrm{N} \otimes_{\mathrm{K}} \mathrm{V}\) 为由 \(\mathbf{V}\) 经标量扩张得到的向量 N-空间；用 \(\tilde{u}\) 表示 \(\mathrm{V}_{(\mathrm{N})}\) 上满足 \(\tilde{u}(x \otimes y) = x \cdot u(y)\) （对 \(x \in \mathbb{N}\) 及 \(y \in \mathbb{V}\) ）的 N-线性型。对每个 \(\sigma \in \Gamma\) ，存在一个 N-线性型 \(h_{\sigma}\) （\(\mathrm{V}_{(\mathrm{N})}\) 上的），使 \(h_{\sigma}(x \otimes y) = x \sigma(y)\) （对 \(x \in \mathbb{N}\) 及 \(y \in \mathbb{V}\) ）。\(\mathrm{V}_{(\mathrm{N})} = \mathrm{N} \otimes_{\mathrm{K}} \mathrm{V}\) 到 \(\mathrm{N} \otimes_{\mathrm{K}} \mathrm{N}\) 中的典范映射是单射。而命题 8 的推论（V，第 64 页）表明线性型 \(h_{\sigma}\) 在 \(\mathrm{V}_{(\mathrm{N})}\) 上的核的交简化为 0。因此（II，第 302 页，推论 1），存在 \(\sigma_1, \ldots, \sigma_n\) （\(\Gamma\) 中的）及 \(a_1, \ldots, a_n\) （\(\mathbb{N}\) 中的）使 \(\tilde{u} = \sum_{i=1}^{n} a_i h_{\sigma_i}\) ，从而 \(u(x) = \sum_{i=1}^{n} a_i \sigma_i(x)\) （对所有 \(x \in V\) ）。

我们给 N 到 N 的全体映射组成的集 \(N^{N}\) 赋以诸因子的离散拓扑的积拓扑。定理 2 意味着，\(\Gamma\) 的元素的、系数在 N 中的线性组合组成的集在 N 到自身的 K-线性映射组成的集中稠密。

定理 3.——设 N 是一个域，\(\Gamma\) 是 N 的一个有限自同构群，K 是 \(\Gamma\) 的不变量域。设 \(n\) 为 \(\Gamma\) 的基数。

\begin{enumerate}
\def\labelenumi{\alph{enumi})}
\item
  我们有 \([\mathbf{N}:\mathbf{K}] = n\) ，且 \(\mathbf{N}\) 是 \(\mathbf{K}\) 的一个伽罗瓦扩张，其伽罗瓦群为 \(\Gamma\) 。
\item
  设 \(\sigma_1, \ldots, \sigma_n\) 是 \(\Gamma\) 的元素，\((x_1, \ldots, x_n)\) 是 \(\mathbf{N}\) 在 \(\mathbf{K}\) 上的一个基，则 \(\det (\sigma_i(x_j)) \neq 0\) 。
\item
  设 \(u\) 是 \(\mathbf{N}\) 到 \(\mathbf{N}\) 中的一个 K-线性映射。存在唯一的族 \((a_{\sigma})_{\sigma \in \Gamma}\) （\(\mathbf{N}\) 的元素组成的），使 \(u(x) = \sum_{\sigma \in \Gamma} a_{\sigma} \sigma(x)\) （对所有 \(x \in \mathbf{N}\) ）。
\end{enumerate}

我们给环 \(N \otimes_{K} N\) 赋以 N-代数结构，其外运算律由 \(\lambda(x \otimes y) = x \otimes \lambda y\) （对 N 中的 \(\lambda, x, y\) ）给出。于是向量 N-空间 \(N \otimes_{K} N\) 的维数是 {[}N: K{]}。积向量 N-空间 \(N^{\Gamma}\) 的维数等于 n。命题 8 的推论（V，第 64 页）中定义的映射 \(\psi\) 是 \(N \otimes_{K} N\) 到 \(N^{\Gamma}\) 上的一个 N-同构，由此 {[}N: K{]} = n。设 \(\Delta\) 为 N 的 K-自同构组成的群。我们有 \(\Gamma \subset \Delta\) ，故 K 是 \(\Delta\) 的不变量域，且 N 是 K 的一个伽罗瓦扩张。进而，由戴德金定理（V，第 27 页，推论 2），\(\Delta\) 的阶至多等于 {[}N: K{]}，而 \(\Gamma\) 的阶等于 {[}N: K{]}，故 \(\Gamma = \Delta\) 。于是 \(\Gamma\) 是 N 在 K 上的伽罗瓦群，这就证明了 a)。

用 \(b)\) 的记号，令 \(f_{i} = \psi(x_{i} \otimes 1)\) ；我们有 \(f_{i}(\sigma) = \sigma(x_{i})\) （对 \(1 \leqslant i \leqslant n\) 及 \(\sigma \in \Gamma\) ）。由于 \(\psi\) 是向量 N-空间的同构，序列 \((f_{1}, \ldots, f_{n})\) 是 \(\mathbf{N}^{\Gamma}\) 在 \(\mathbf{N}\) 上的一个基，由此 \(\det(f_{j}(\sigma_{i})) \neq 0\) ，即

\[
\det \left(\sigma_ {i} (x _ {j})\right) \neq 0.
\]

这就证明了 \(b\) )。

最后，\(c)\) 由定理 2（V，第 65 页）得出，该定理证明了族 \((a_{\sigma})_{\sigma \in \Gamma}\) 的存在性，使 \(u(x) = \sum_{\sigma \in \Gamma} a_{\sigma} \sigma(x)\) （对所有 \(x \in \mathbf{N}\) ）；并由戴德金定理

（V，第 27 页，推论 2）证明了 \((a_{\sigma})_{\sigma \in \Gamma}\) 的唯一性。

